% 第8章 数组和算法基础（教材 8.1--8.7）
% 只写内容：不重复 \usepackage，不改字号与页面设置

\section{一维数组的定义和初始化}
数组（Array）是\textbf{相同类型}数据的有序集合，在内存中占据一段\textbf{连续}的存储单元。数组把"一批同类型数据"打包成一个整体，用一个名字统一管理，再用下标区分其中的每个元素，这正是后续排序、查找、字符串、矩阵运算的基础。

\subsection{数组的定义}
定义一维数组的一般形式为：
\begin{verbatim}
类型关键字  数组名[常量表达式];
\end{verbatim}
例如 \texttt{int score[40];} 定义了一个名为 \texttt{score} 的整型数组，可存放 40 个 \texttt{int} 型数据。

语法规则有四点：
\begin{enumerate}
  \item \textbf{长度必须是常量表达式}。标准 C（C89/C90）要求方括号内是编译期可求值的整型常量表达式，因此 \texttt{int n=40; int a[n];} 在 C89 下是错的；C99 起允许变长数组（VLA），但教材与多数评测环境仍按"常量"要求。
  \item \textbf{下标从 0 开始}，合法下标范围是 $0\sim n-1$。长度为 $n$ 的数组，最后一个元素是 \texttt{a[n-1]}，\texttt{a[n]} 已经越界。
  \item \textbf{数组名是这批存储单元的首地址}，是一个地址常量，不能被赋值、不能做自增自减：\texttt{a = b;} 与 \texttt{a++;} 都是非法表达式。
  \item \textbf{编译器不做下标越界检查}。越界访问在语法上完全合法，运行时可能悄悄改掉别的变量，也可能直接崩溃——这是 C 语言最典型的"沉默错误"。
\end{enumerate}

\subsection{数组的初始化}
数组可以在定义时初始化，常见形式见下表。

\begin{table}[htbp]
\centering
\caption{一维数组初始化形式对照}
\footnotesize
\begin{tabularx}{\textwidth}{@{}p{4.6cm}Xp{2.6cm}@{}}
\toprule
初始化写法 & 含义 & 未写出元素的值 \\
\midrule
\texttt{int a[5] = \{1,2,3,4,5\};} & 全部元素依次赋初值 & —— \\
\texttt{int a[5] = \{1,2,3\};} & 部分初始化 & 其余自动为 0 \\
\texttt{int a[5] = \{0\};} & 常见的"整体清零"写法 & 全部为 0 \\
\texttt{int a[] = \{1,2,3,4,5\};} & 省略长度，由初值个数决定 & 长度为 5 \\
\texttt{int a[5];} & 只定义不初始化 & 不确定（随机残留值） \\
\bottomrule
\end{tabularx}
\end{table}

\corepoint{部分初始化时"补 0"只发生在\textbf{定义时的初始化}里；若是先定义 \texttt{int a[5];} 再在函数内逐个赋值，没赋到的元素仍是随机值。想把一个已存在的数组清零，应使用循环或 \texttt{memset}，不能写 \texttt{a = \{0\};}——数组名不可整体赋值。}

\subsection{内存模型与元素访问}
\texttt{int score[40];} 在 32 位环境下的内存布局可以想成 40 个紧挨着的抽屉，每个抽屉 4 字节，总长 $40\times4=160$ 字节。第 $i$ 个元素的首地址满足
\keyformula{\&\texttt{score}[i] = \&\texttt{score}[0] + i \times \texttt{sizeof(int)}}
"连续存放"带来两个直接后果：一是可以用"首地址 + 偏移量"线性地走遍整个数组；二是下标运算的本质就是地址运算，这一点在第 11 章讲指针与数组的关系时会彻底展开。

遍历数组的标准写法是"下标从 0 到 n-1 的 for 循环"：
\begin{lstlisting}
int a[10], i, sum = 0;
for (i = 0; i < 10; i++) {   /* 注意是 < 10，不是 <= 10 */
    scanf("%d", &a[i]);
    sum += a[i];
}
printf("average = %.2f\n", (double)sum / 10);
\end{lstlisting}

\section{二维数组的定义和初始化}
\subsection{二维数组的定义}
定义形式为 \texttt{类型 数组名[行数][列数];}，例如 \texttt{int score[3][4];} 表示 3 行 4 列共 12 个 \texttt{int} 元素。元素用两个下标引用：\texttt{score[i][j]}，其中 $0\le i\le2$，$0\le j\le3$。

\textbf{二维数组在内存中并不是"二维"的}：C 语言按\textbf{行优先}（row-major）顺序把二维数组拉直成一维存放，先存第 0 行全部元素，再存第 1 行，依此类推。因此 \texttt{score[1][0]} 紧跟在 \texttt{score[0][3]} 之后。设列数为 $N$，则
\keyformula{\&\texttt{a}[i][j] = \texttt{a} + (i \times N + j)\times \texttt{sizeof(元素类型)}}

\begin{table}[htbp]
\centering
\caption{\texttt{int a[3][4]} 的行优先存储布局}
\footnotesize
\begin{tabular}{@{}llllll@{}}
\toprule
内存顺序 & 0 & 1 & 2 & $\cdots$ & 11 \\
\midrule
元素 & \texttt{a[0][0]} & \texttt{a[0][1]} & \texttt{a[0][2]} & $\cdots$ & \texttt{a[2][3]} \\
偏移（元素个数） & 0 & 1 & 2 & $\cdots$ & 11 \\
\bottomrule
\end{tabular}
\end{table}

\subsection{二维数组的初始化}
\begin{itemize}
  \item \textbf{分行初始化}（推荐，可读性最好）：
\begin{lstlisting}
int a[3][4] = {{1, 2, 3, 4},
               {5, 6, 7, 8},
               {9, 10, 11, 12}};
\end{lstlisting}
  \item \textbf{连续初始化}：\texttt{int a[3][4] = \{1,2,3,4,5,6,7,8,9,10,11,12\};}，编译器按行优先自动"折行"。
  \item \textbf{部分初始化}：\texttt{int a[3][4] = \{\{1\},\{5,6\}\};}，未写出的元素全部补 0。
  \item \textbf{省略行数}：\texttt{int a[][4] = \{\{1,2,3,4\},\{5,6,7,8\}\};} 合法；但\textbf{列数永远不能省略}，因为编译器要靠列数计算"一行有多长"才能确定行数。
\end{itemize}

\subsection{二维数组的遍历}
遍历二维数组要用双重循环，通常\textbf{外层控制行、内层控制列}：
\begin{lstlisting}
#define M 3
#define N 4
int a[M][N], i, j;
for (i = 0; i < M; i++) {
    for (j = 0; j < N; j++)
        printf("%4d", a[i][j]);
    printf("\n");            /* 一行输出完再换行 */
}
\end{lstlisting}

\corepoint{若把两层循环写反（外层列、内层行），结果虽然也能"算对"，但访问顺序与内存顺序相反，缓存命中率变差；更常见的错误是把循环边界写成行数/列数互换，导致越界。养成"M 行 N 列、外 i 内 j"的固定习惯可以避免大半错误。}

\section{向函数传递一维数组}
\subsection{形参的三种等价写法}
把一维数组传给函数时，下面三种形参写法完全等价：
\begin{lstlisting}
void PrintArray(int a[10], int n);   /* 方括号里的 10 只是"提示"，不起约束作用 */
void PrintArray(int a[],   int n);
void PrintArray(int *a,    int n);   /* 编译器实际看到的就是这个 */
\end{lstlisting}
原因是：\textbf{数组名作实参时，传递的是首元素地址，不是整个数组}。C 语言没有"按值传递整个数组"这回事，形参数组在函数内部已经退化为指针，所以
\begin{itemize}
  \item 方括号内写不写长度都一样，写了也不检查；
  \item 函数无法从形参得知数组有多少个元素，必须\textbf{另外传一个长度参数} $n$；
  \item 在函数内 \texttt{sizeof(a)} 得到的是\textbf{指针的大小}（4 或 8 字节），不是数组总字节数。
\end{itemize}

\subsection{地址传递的效果}
因为传的是地址，函数内对 \texttt{a[i]} 的修改会直接作用在实参数组上——这与普通变量"按值传递、函数内改了外面不变"的行为截然不同。

\begin{table}[htbp]
\centering
\caption{普通变量与数组作函数参数的区别}
\footnotesize
\begin{tabularx}{\textwidth}{@{}p{2.6cm}p{4.2cm}X@{}}
\toprule
 & 普通变量 & 一维数组 \\
\midrule
实参形式 & \texttt{f(x)} & \texttt{f(a, n)} \\
传递内容 & 值的一份副本 & 首元素地址 \\
函数内修改 & 不影响实参 & 直接影响实参元素 \\
函数内 \texttt{sizeof} & 变量本身的大小 & 指针的大小 \\
需要额外传长度 & 不需要 & 必须传 \\
\bottomrule
\end{tabularx}
\end{table}

\begin{lstlisting}
/* 求数组元素最大值：形参用 const 表明"只读"，是一种防御式编程习惯 */
int FindMax(const int a[], int n)
{
    int i, max = a[0];
    for (i = 1; i < n; i++)
        if (a[i] > max) max = a[i];
    return max;
}
\end{lstlisting}

\section{排序和查找}
在数组中搜索一个特定元素的过程称为\textbf{查找}（search），把数组元素按关键字重新排列的过程称为\textbf{排序}（sort）。被查找的值称为\textbf{查找键}（searchKey）。排序与查找是数组最重要的两类应用，也是"算法"这一概念的入口。

\subsection{交换法排序}
思路：第 $i$ 轮把 \texttt{a[i]} 与其后所有元素逐个比较，\textbf{一旦发现更大（或更小）就立即交换}。$n$ 个元素共需 $n-1$ 轮。
\begin{lstlisting}
void ExchangeSort(int a[], int n)      /* 由大到小 */
{
    int i, j, temp;
    for (i = 0; i < n - 1; i++)
        for (j = i + 1; j < n; j++)
            if (a[j] > a[i]) {         /* 边比较边交换，交换次数多 */
                temp = a[i]; a[i] = a[j]; a[j] = temp;
            }
}
\end{lstlisting}
\textbf{特点}：比较次数固定为 $n(n-1)/2$，但交换次数多，最坏情况下每比较一次都要搬三次数据。

\subsection{冒泡法排序}
思路：相邻两个元素比较，逆序就交换。每一轮结束后，当前未排序部分中的最值会"冒"到一端，因此下一轮可以少比一个。
\begin{lstlisting}
void BubbleSort(int a[], int n)        /* 由小到大 */
{
    int i, j, temp, flag;
    for (i = 0; i < n - 1; i++) {
        flag = 0;
        for (j = 0; j < n - 1 - i; j++)   /* 已就位的尾部不再比较 */
            if (a[j] > a[j + 1]) {
                temp = a[j]; a[j] = a[j + 1]; a[j + 1] = temp;
                flag = 1;
            }
        if (flag == 0) break;             /* 本轮无交换，已经有序，提前结束 */
    }
}
\end{lstlisting}
\textbf{特点}：最好情况（原本有序）只需 $n-1$ 次比较、0 次交换；加 \texttt{flag} 优化后可以提前退出。

\subsection{选择法排序}
思路：第 $i$ 轮先在 \texttt{a[i..n-1]} 中\textbf{记下最值的下标}，本轮结束再交换一次。与交换法排序的区别只有一点：\textbf{交换发生在每轮末尾，而不是每次比较之后}。
\begin{lstlisting}
void SelectionSort(int a[], int n)     /* 由小到大 */
{
    int i, j, k, temp;
    for (i = 0; i < n - 1; i++) {
        k = i;                          /* k 记录本轮最小元素的下标 */
        for (j = i + 1; j < n; j++)
            if (a[j] < a[k]) k = j;
        if (k != i) {                   /* 每轮至多交换一次 */
            temp = a[i]; a[i] = a[k]; a[k] = temp;
        }
    }
}
\end{lstlisting}

\begin{table}[htbp]
\centering
\caption{三种简单排序算法对照（$n$ 个元素）}
\footnotesize
\begin{tabularx}{\textwidth}{@{}p{2.3cm}p{2.2cm}p{2.3cm}Xp{2.1cm}@{}}
\toprule
算法 & 比较次数 & 交换次数 & 关键特征 & 稳定性 \\
\midrule
交换法排序 & 约 $n^2/2$ & 最坏约 $n^2/2$ & 每次比较都可能交换 & 不稳定 \\
冒泡法排序 & 最坏 $n^2/2$ & 逆序对个数 & 相邻比较、最值逐步"冒"出 & 稳定 \\
选择法排序 & 约 $n^2/2$ & 至多 $n-1$ & 先记下标、轮末交换一次 & 不稳定 \\
\bottomrule
\end{tabularx}
\end{table}

\subsection{线性查找（顺序查找）}
思路最直白：从数组一端开始，用查找键 \texttt{key} 依次与每个元素比较，相等即返回下标；比完全部仍未相等，则返回 $-1$ 表示"未找到"。
\begin{lstlisting}
int LinSearch(const int a[], int n, int key)
{
    int i;
    for (i = 0; i < n; i++)
        if (a[i] == key) return i;   /* 找到，返回下标 */
    return -1;                       /* 全部比完仍未找到 */
}
\end{lstlisting}
\textbf{特点}：对数据\textbf{没有任何前提要求}（可以无序），实现最简单；代价是效率低——最坏要比较 $n$ 次，平均约 $(n+1)/2$ 次，时间复杂度 $O(n)$。数据无序时，线性查找是唯一可行的办法。

\subsection{折半查找}
\textbf{前提条件：数组元素已按关键字有序排列}。在此前提下，折半查找（binary search，也称二分查找、对分搜索）每次比较都能排除一半元素。

算法步骤：
\begin{enumerate}
  \item 设查找区间为 \texttt{[low, high]}，初值 \texttt{low = 0}，\texttt{high = n - 1}；
  \item 取中间位置 \texttt{mid = (low + high) / 2}，把 \texttt{a[mid]} 与查找键 \texttt{key} 比较；
  \item 若相等，查找成功，返回 \texttt{mid}；
  \item 若 \texttt{key < a[mid]}（升序表），说明目标在前半段，令 \texttt{high = mid - 1}；否则令 \texttt{low = mid + 1}；
  \item 重复 2--4，直到 \texttt{low > high}，此时区间为空，查找失败。
\end{enumerate}
\keyformula{\texttt{mid} = \left\lfloor \frac{\texttt{low}+\texttt{high}}{2} \right\rfloor}

\begin{lstlisting}
int BinarySearch(const int a[], int n, int key)
{
    int low = 0, high = n - 1, mid;
    while (low <= high) {              /* 区间非空才继续 */
        mid = (low + high) / 2;
        if (key == a[mid]) return mid; /* 找到，返回下标 */
        else if (key < a[mid]) high = mid - 1;
        else                   low  = mid + 1;
    }
    return -1;                         /* 未找到 */
}
\end{lstlisting}

效率：每比较一次就排除一半元素，因此最多所需的比较次数约为 $\log_2 n + 1$ 次。$n=1024=2^{10}$ 时最坏只要 10 次左右比较，而顺序查找最坏要 1024 次。这就是"有序"换来的效率。

\corepoint{折半查找三处易错：（1）\textbf{忘记先排序}——对无序数组折半，结果无意义；（2）循环条件写成 \texttt{low < high}，会漏掉区间里只剩一个元素的情况；（3）区间收缩写成 \texttt{high = mid} 或 \texttt{low = mid}（没有 $\pm1$），可能造成死循环。}

\begin{table}[htbp]
\centering
\caption{顺序查找与折半查找对照}
\footnotesize
\begin{tabularx}{\textwidth}{@{}p{2.4cm}p{2.4cm}Xp{2.6cm}@{}}
\toprule
算法 & 对数据的要求 & 基本思想 & 平均效率 \\
\midrule
顺序查找（线性查找） & 无要求，可无序 & 从一端逐个与查找键比较 & $O(n)$ \\
折半查找（二分查找） & \textbf{必须有序} & 每次与中间元素比较，排除一半 & $O(\log_2 n)$ \\
\bottomrule
\end{tabularx}
\end{table}

\section{向函数传递二维数组}
二维数组作函数参数时，形参必须写成"数组的数组"，\textbf{第二维（列数）不能省略}：
\begin{lstlisting}
void PrintMatrix(int a[][4], int rows);      /* 正确：给出列数 4 */
void PrintMatrix(int (*a)[4], int rows);     /* 与上一行完全等价 */
/* void PrintMatrix(int a[][], int rows);     错误：列数未知，无法计算地址 */
\end{lstlisting}
原因是：函数内部要用 \texttt{a[i][j]} 定位元素，必须先知道"一行有多少个元素"才能算出偏移量 $i\times N+j$；行数只影响循环边界，所以可以作为普通参数传入。

\begin{lstlisting}
/* 求每行元素的平均值：rows 为实际行数，列数固定为 N */
#define N 4
void RowAverage(int a[][N], int rows)
{
    int i, j, sum;
    for (i = 0; i < rows; i++) {
        sum = 0;
        for (j = 0; j < N; j++) sum += a[i][j];
        printf("row %d average = %.2f\n", i, (double)sum / N);
    }
}
\end{lstlisting}

\corepoint{不要以为 \texttt{int **} 能接 \texttt{int a[3][4]}。二维数组是一片连续内存，\texttt{a} 的类型是"指向含 4 个 int 的数组"的指针，与"指向指针的指针" \texttt{int **} 是两种完全不同的类型；混用会在运行期按错误的地址取值（详见第 11 章）。}

\section{本章知识点小结}
\begin{itemize}
  \item 数组是\textbf{同类型、连续存储}的元素集合；下标从 0 起，合法范围 $0\sim n-1$。
  \item 定义时长度须为常量表达式；初始化分全部、部分（自动补 0）、省略长度三种；数组名是地址常量，不能整体赋值、不能自增。
  \item 一维数组元素地址 $\&\texttt{a}[i]=\&\texttt{a}[0]+i\times\texttt{sizeof(元素)}$；二维数组按\textbf{行优先}拉直存放，$\&\texttt{a}[i][j]=\texttt{a}+(iN+j)\times\texttt{sizeof(元素)}$。
  \item 数组作函数实参传递的是首地址，形参 \texttt{a[]}、\texttt{a[10]}、\texttt{*a} 三者等价，必须另传长度；函数内 \texttt{sizeof} 失效。
  \item 二维数组作形参时列数不可省；行数用普通参数传入。
  \item 排序：交换法（边比边换）、冒泡法（相邻比较、可提前退出）、选择法（记下标、轮末交换）。查找：顺序查找无前提、$O(n)$；折半查找须有序、$O(\log_2 n)$。
\end{itemize}

\section{本章常见错误小结}
\begin{enumerate}
  \item \textbf{用圆括号引用数组元素}：写成 \texttt{a(3)} 或 \texttt{a[2,3]}（把行列写进同一个方括号）。前者被当成函数调用，后者是逗号表达式的值 3，都不是"第 3 个元素"。
  \item \textbf{下标越界}：\texttt{for (i = 1; i <= n; i++) a[i] = 0;} 访问了 \texttt{a[n]}，编译器不报错，却已破坏相邻内存。
  \item \textbf{给数组名赋值}：\texttt{a = b;}、\texttt{a++;}、"数组整体赋初值" \texttt{a = \{0\};} 均非法；数组整体拷贝要逐元素或用 \texttt{memcpy}。
  \item \textbf{函数内用 sizeof 求数组长度}：\texttt{int len = sizeof(a) / sizeof(a[0]);} 在形参 \texttt{a} 是数组名时可用，在函数内 \texttt{a} 已退化为指针，结果恒为 1 或 2，是经典陷阱。
  \item \textbf{二维数组形参省掉列数}：\texttt{void f(int a[][])} 编译不过；即使写成 \texttt{int **a}，与二维数组也不是同一类型。
  \item \textbf{对无序数组折半查找}：前提条件不成立，结果错误；应先排序。
  \item \textbf{行列循环写反或边界用错}：遍历 \texttt{a[3][4]} 时用 \texttt{j < 3} 控制列，少访问了一列。
\end{enumerate}
